\documentclass[12pt, a4paper]{article}

% --- Packages ---
\usepackage[utf8]{inputenc}
\usepackage[T1]{fontenc}
\usepackage{geometry}
\geometry{left=2.5cm, right=2.5cm, top=2.5cm, bottom=2.5cm}
\usepackage{times}
\usepackage[numbers,sort&compress]{natbib}
\usepackage{hyperref}
\usepackage{graphicx}
\usepackage{titlesec}
\usepackage{enumitem}
\usepackage{setspace}
\usepackage{amsmath}
\usepackage{amssymb}
\usepackage{booktabs}
\usepackage{array}
\usepackage{multirow}
\usepackage{xcolor}

\hypersetup{
    colorlinks=true,
    linkcolor=blue,
    filecolor=magenta,
    urlcolor=cyan,
    citecolor=blue,
    pdfencoding=auto
}

\onehalfspacing
\setlength{\parskip}{0.5em}
\titlespacing{\section}{0pt}{14pt}{6pt}
\titlespacing{\subsection}{0pt}{10pt}{4pt}
\titlespacing{\subsubsection}{0pt}{8pt}{3pt}

% =========================================================
\title{\textbf{Multi-Agent Reinforcement Learning for Autonomous Driving:\\
Algorithms, Coordination, and Emerging Paradigms}\\[0.4em]
{\large \textit{A Comprehensive Survey}}}
\author{\large Literature Review}
\date{\today}

\begin{document}

\maketitle

% =========================================================
\begin{abstract}
Autonomous driving (AD) is one of the most demanding challenges in modern
artificial intelligence, requiring real-time perception, intent prediction of
heterogeneous road users, and safe, socially compliant decision-making in a
continuously evolving multi-agent environment.
This survey provides a systematic and comprehensive review of the field,
organized around a three-layer taxonomy that traces the evolution from
single-agent foundations to fully cognitive driving systems.

\textbf{Layer 1---Foundational Single-Agent RL.}
We begin by formalizing the AD decision problem as a Partially Observable
Markov Decision Process (POMDP)~\cite{hubmann2018automated} and surveying the
algorithmic building blocks that underpin all subsequent work: end-to-end
behavior cloning~\cite{bojarski2016end}, augmentation strategies to overcome
covariate shift~\cite{bansal2018chauffeur}, conditional imitation
learning~\cite{codevilla2018end}, and policy gradient optimization including
PPO~\cite{schulman2017proximal} and its deep RL derivatives~\cite{sallab2017deep,kiran2021deep}.

\textbf{Layer 2---MARL Algorithms and Coordination Protocols.}
We then extend the formalism to the Partially Observable Stochastic Game (POSG)
framework~\cite{littman1994markov,bernstein2002complexity}, where the
non-stationarity and exponential joint-action space of traffic demand specialized
multi-agent methods.  We survey value-decomposition approaches
(VDN~\cite{sunehag2018value}, QMIX~\cite{rashid2018qmix}, QTRAN~\cite{son2019qtran}),
actor-critic CTDE methods (MADDPG~\cite{lowe2017multi},
COMA~\cite{foerster2018counterfactual}, MAPPO~\cite{yu2021surprising},
HATRPO~\cite{kuba2021hatrpo}), role-based architectures
(RODE~\cite{wang2021rode}, ROMA~\cite{wang2020roma}), hierarchical
MARL~\cite{vezhnevets2017feudal,chu2020multi}, and game-theoretic
approaches~\cite{shapley1953stochastic,littman1994markov,hu2003nash,foerster2018lola}.
For communication and coordination, we trace the progression from early
learned protocols~\cite{foerster2016learning,sukhbaatar2016learning} to
targeted attention-based messaging~\cite{das2019tarmac,jiang2018learning},
graph-structured communication~\cite{jiang2019graph,kim2019learning}, and
V2X cooperative perception~\cite{wang2020v2vnet}.
On the representation side, we analyze the evolution from grid social
pooling~\cite{alahi2016social} to vectorized GNNs~\cite{gao2020vectornet,liang2020lanegcn},
multi-modal prediction~\cite{gupta2018social,salzmann2020trajectron,zhao2020tnt},
and bidirectional prediction--planning coupling~\cite{song2020pip}.

\textbf{Layer 3---Scalable and Transformer/LLM-Augmented Agents.}
The final layer surveys the nascent paradigm of cognitive driving agents,
including Multi-Agent Transformer architectures~\cite{wen2022multi},
Decision Transformers~\cite{chen2021decision}, and LLM-based systems
that inject common-sense reasoning and few-shot
generalization~\cite{fu2023drive,wen2023dilu,mao2023gpt,sha2023languagempc}.
We also survey critical simulation environments (CARLA~\cite{dosovitskiy2017carlaopenurbandriving},
SMARTS~\cite{zhou2020smarts}), benchmarks
(NuPlan~\cite{caesar2021nuplan}, Argoverse~\cite{chang2019argoverse},
Waymo~\cite{ettinger2021large}), and safety frameworks~\cite{shalev2016safe,achiam2017constrained}.

Across all paradigms, we synthesize the key technical contributions,
identify the precise limitations each work addresses, and draw cross-cutting
connections.  The survey covers 100 primary references spanning 2003 to 2024,
and concludes with a structured analysis of open challenges: formal safety
verification under distributional shift, bandwidth-constrained V2X communication,
real-time LLM inference, and the need for generative world models to cover
long-tail corner cases.
\end{abstract}

\noindent\textbf{Keywords:} Autonomous Driving, Multi-Agent Reinforcement Learning,
Centralized Training Decentralized Execution, Trajectory Prediction, Graph Neural
Networks, Vehicle-to-Vehicle Communication, Large Language Models, Cognitive Agents,
Safety, Simulation.

\newpage
\tableofcontents
\newpage

% ==========================================================
\section{Introduction}
\label{sec:intro}

The pursuit of fully autonomous driving (AD) represents one of the most
transformative yet technically demanding frontiers in modern artificial
intelligence.  It requires agents to simultaneously perceive complex,
dynamic environments with high fidelity; predict the intentions of
heterogeneous road users ranging from aggressive merging vehicles to
unpredictable pedestrians; and execute safe, efficient, and socially
compliant maneuvers under continuous time pressure.  Critically, these
tasks are not independent: the optimal action for the ego-vehicle depends
on the predicted behavior of surrounding agents, which in turn depends on
the ego-vehicle's own future actions.  This mutual dependence renders
autonomous driving a canonical \emph{multi-agent} decision problem in
which the performance of any single policy is fundamentally coupled to the
policies of all other actors in the scene.

Despite decades of engineering effort, achieving the level of robustness and
generalization required for unrestricted public deployment on diverse road
types and in novel scenarios remains an open research problem.  Failures
continue to occur precisely in situations that are statistically rare (so-called
\emph{long-tail} events) yet critically important for safety: unusual road
geometries, ambiguous right-of-way negotiations, adversarial pedestrian
behavior, or unexpected sensor degradation.  The historical arc of the field
can be understood as a sequence of three successive paradigm shifts, each
motivated by well-defined limitations of its predecessor.

% ------------------------------------------------------------------
\subsection{Motivation and Problem Statement}
\label{sec:motivation}

Classical autonomous driving systems relied on a modular pipeline
architecture that decomposes the complex driving task into distinct,
hierarchical sub-systems: perception, localization, prediction, planning,
and control~\cite{caesar2021nuplan}.  While this ``divide-and-conquer''
approach offers engineering interpretability and allows individual modules
to be validated in isolation, it suffers from systematic limitations.
Most importantly, it is prone to \emph{cascading errors}: a minor
imperfection in the perception module (e.g., a flickering bounding box at
the boundary of sensor range) can propagate downstream, destabilizing the
planner and leading to conservative, suboptimal, or even dangerous
behavior.  Furthermore, hand-crafted rules in the planning module fail to
capture the nuanced, interactive nature of human driving, resulting in
behavior that is either overly conservative or strategically naive.

The promise of data-driven, end-to-end learning as an alternative was first
demonstrated at scale by \citeauthor{bojarski2016end}~\cite{bojarski2016end},
who showed that a convolutional neural network mapping raw front-camera images
directly to steering commands could navigate real roads without explicit
intermediate representations.  However, the limitations of pure imitation
learning quickly became apparent: the \emph{covariate shift} problem, wherein
the agent drifts into states absent from the training distribution and has no
learned recovery mechanism, remained a fundamental obstacle~\cite{bansal2018chauffeur}.
More deeply, pure imitation cannot model the \emph{strategic interactions}
between agents: the response of a merging vehicle depends not only on its own
preferences but also on its prediction of the ego-vehicle's intent, creating a
feedback loop that single-agent behavior cloning cannot represent.

This interactive structure formally defines autonomous driving as a
\emph{multi-agent} problem.  In a traffic environment with $N$ vehicles,
each agent $i$ observes a partial view $o^i_t$ of the world state $s_t$,
executes an action $a^i_t$ from its action space $\mathcal{A}^i$, and
receives a reward $R^i(s_t, \mathbf{a}_t)$ that depends on the
\emph{joint} action vector $\mathbf{a}_t = (a^1_t, \ldots, a^N_t)$.
The transition dynamics
$T(s_{t+1} \mid s_t, \mathbf{a}_t)$ couple all agents together, so
no single agent can treat the environment as stationary while other
agents are simultaneously adapting their
policies~\cite{zhang2021multi,busoniu2008comprehensive}.

This non-stationarity violates the foundational Markov assumption required
for convergence of standard single-agent RL algorithms.  \citeauthor{shalev2016safe}~\cite{shalev2016safe}
further show that even if non-stationarity is managed, the gradient
variance of policy gradient methods scales as $O(1/p^2)$ for accident
probability $p \ll 1$, making unconstrained RL impractically slow for
safety-critical applications.  The central research question motivating
this survey is therefore:

\begin{quote}
\emph{How have MARL algorithms, coordination protocols, and cognitive
architectures evolved to address the cooperative and competitive challenges
inherent in multi-agent traffic environments, and what are the remaining
open problems?}
\end{quote}

% ------------------------------------------------------------------
\subsection{Driving Forces Behind Paradigm Shifts}
\label{sec:drivers}

The field has not evolved uniformly; rather, it has undergone discrete
paradigm shifts each driven by a convergence of factors.  We identify four
primary categories of driving forces:

\subsubsection*{(1) Algorithmic Breakthroughs}
The development of deep Q-networks (DQN) by \citeauthor{mnih2015human}~\cite{mnih2015human}
demonstrated that deep neural networks could learn superhuman strategies for
Atari games directly from raw pixels, immediately suggesting that similar
methods could be applied to driving perception and control.  Policy gradient
methods (REINFORCE, A3C/A2C)~\cite{mnih2016asynchronous} and their proximal
trust-region variants (TRPO, PPO~\cite{schulman2017proximal}) provided stable
and scalable optimization for continuous-action driving tasks.  The discovery
of the CTDE (Centralized Training, Decentralized Execution) framework by
\citeauthor{kraemer2016multi}~\cite{kraemer2016multi} and its instantiations
in MADDPG~\cite{lowe2017multi} and QMIX~\cite{rashid2018qmix} unlocked
principled cooperative MARL.  Graph neural networks
(GAT~\cite{velickovic2018graph}, GCN-based policy networks~\cite{jiang2019graph})
provided the relational inductive biases needed to encode agent-to-agent
and agent-to-map relationships compactly.  Most recently, Transformer
architectures~\cite{vaswani2017attention} and large language
models~\cite{brown2020language} have begun reshaping the cognitive layer of
driving agents~\cite{wen2022multi,fu2023drive,sha2023languagempc}.

\subsubsection*{(2) Hardware and Compute Scaling}
The availability of large GPU clusters has made it practical to train policies
on billions of environment steps, a prerequisite for sample-hungry MARL
algorithms.  Real-world sensor packages---combining high-resolution LiDAR,
multi-camera arrays, GPS/IMU, and radar---have reached cost levels that enable
fleet-scale data collection, providing the large-scale supervision necessary for
pre-training foundation models for driving~\cite{caesar2020nuscenes}.

\subsubsection*{(3) Dataset and Simulation Ecosystem Maturation}
High-fidelity, open-source driving simulators (CARLA~\cite{dosovitskiy2017carlaopenurbandriving},
SUMO~\cite{krajzewicz2012recent}) and domain-specific MARL training
platforms (SMARTS~\cite{zhou2020smarts}, MACAD-Gym~\cite{palanisamy2020multi},
Flow~\cite{wu2017flow}) have dramatically reduced the barrier to multi-agent
experimentation.  Simultaneously, large-scale real-world driving datasets
(nuScenes~\cite{caesar2020nuscenes}, Argoverse~\cite{chang2019argoverse},
Waymo Open Motion~\cite{ettinger2021large}) and closed-loop planning
benchmarks (NuPlan~\cite{caesar2021nuplan}) have standardized evaluation
and exposed the limitations of open-loop metrics.

\subsubsection*{(4) Real-World Deployment Demands}
Commercial robotaxi programs (Waymo, Cruise, Baidu Apollo) have pushed
research toward long-tail robustness, edge-case handling, and
interpretability---precisely the domains where purely data-driven methods
struggle.  Emerging V2X infrastructure standards (C-V2X, DSRC) have
created new opportunities for cooperative perception and planning that
transcend what any single-vehicle sensor suite can
achieve~\cite{wang2020v2vnet}.

% ------------------------------------------------------------------
\subsection{Survey Scope and Literature Selection Methodology}
\label{sec:scope}

\subsubsection*{Topic Coverage}
This survey covers the following interconnected topics:
\begin{itemize}[noitemsep]
  \item \textbf{Formal foundations}: POMDP, Markov games~\cite{littman1994markov},
        stochastic games~\cite{shapley1953stochastic},
        Dec-POMDP~\cite{bernstein2002complexity};
  \item \textbf{Single-agent RL for driving}: behavior cloning~\cite{bojarski2016end},
        covariate shift mitigation~\cite{bansal2018chauffeur},
        conditional imitation~\cite{codevilla2018end},
        risk-sensitive RL~\cite{kahn2017uncertainty},
        hierarchical RL~\cite{chen2019model};
  \item \textbf{MARL algorithms}: value decomposition
        (VDN~\cite{sunehag2018value}, QMIX~\cite{rashid2018qmix},
        QTRAN~\cite{son2019qtran}, MAVEN~\cite{mahajan2019maven}),
        actor-critic (MADDPG~\cite{lowe2017multi},
        COMA~\cite{foerster2018counterfactual},
        MAPPO~\cite{yu2021surprising}, HATRPO~\cite{kuba2021hatrpo},
        FACMAC~\cite{peng2021facmac}),
        role-based (RODE~\cite{wang2021rode}, ROMA~\cite{wang2020roma}),
        hierarchical~\cite{vezhnevets2017feudal,chu2020multi},
        game-theoretic (Nash Q-learning~\cite{hu2003nash},
        LOLA~\cite{foerster2018lola}, PSRO~\cite{lanctot2017unified});
  \item \textbf{Communication and coordination}: DIAL/RIAL~\cite{foerster2016learning},
        CommNet~\cite{sukhbaatar2016learning},
        TarMAC~\cite{das2019tarmac},
        ATOC~\cite{jiang2018learning},
        IC3Net~\cite{singh2019learning},
        structured communication~\cite{kim2019learning},
        message dropout~\cite{kim2019message},
        bandwidth-constrained protocols~\cite{wang2020bandwidth},
        VBC~\cite{wang2021vbc},
        V2X cooperative perception~\cite{wang2020v2vnet};
  \item \textbf{Interaction-aware representation and prediction}: Social
        LSTM~\cite{alahi2016social}, Social GAN~\cite{gupta2018social},
        Trajectron++~\cite{salzmann2020trajectron},
        VectorNet~\cite{gao2020vectornet}, LaneGCN~\cite{liang2020lanegcn},
        TNT~\cite{zhao2020tnt}, MultiPath~\cite{chai2019multipath},
        PiP~\cite{song2020pip};
  \item \textbf{Multi-agent autonomous driving applications}: traffic signal
        control~\cite{wei2019colight,wei2019presslight,zheng2019learning},
        cooperative lane change~\cite{wang2018urban},
        highway driving~\cite{isele2018navigating},
        mixed-autonomy traffic~\cite{wu2017flow,stern2018dissipation},
        POMDP planning~\cite{hubmann2018automated},
        game-theoretic driving~\cite{fisac2019hierarchical,sadigh2016planning,fridovich2020efficient},
        intent-aware planning~\cite{schwarting2019social,schwarting2021deep};
  \item \textbf{Transformer/LLM-based agents}: Multi-Agent
        Transformer~\cite{wen2022multi}, Decision Transformer~\cite{chen2021decision},
        GTrXL~\cite{parisotto2020stabilizing},
        DiLu~\cite{wen2023dilu},
        GPT-Driver~\cite{mao2023gpt},
        Drive Like a Human~\cite{fu2023drive},
        LanguageMPC~\cite{sha2023languagempc};
  \item \textbf{Safety and robustness}: constrained policy
        optimization~\cite{achiam2017constrained}, safe RL via
        shielding~\cite{alshiekh2018safe}, risk-aware
        MARL~\cite{prashanth2016variance}, responsibility-sensitive
        safety~\cite{shalev2016safe};
  \item \textbf{Sim-to-real transfer}: domain
        randomization~\cite{tobin2017domain}, meta-RL (MAML)~\cite{finn2017model};
  \item \textbf{Simulation and benchmarks}: CARLA~\cite{dosovitskiy2017carlaopenurbandriving},
        SUMO~\cite{krajzewicz2012recent}, SMARTS~\cite{zhou2020smarts},
        SMAC~\cite{samvelyan2019starcraft},
        nuScenes~\cite{caesar2020nuscenes}, NuPlan~\cite{caesar2021nuplan},
        Argoverse~\cite{chang2019argoverse},
        Waymo Open Motion~\cite{ettinger2021large}.
\end{itemize}

\subsubsection*{Time Span and Search Strategy}
The survey spans foundational theoretical work from 1953
(Shapley's stochastic games~\cite{shapley1953stochastic}) through the
present (2024), with the main body focused on the period 2016--2024 during
which deep learning methods came to dominate the field.  The literature was
assembled through:
\begin{enumerate}[noitemsep]
  \item Systematic keyword searches on IEEE Xplore, ACM Digital Library,
        Semantic Scholar, and arXiv using the terms
        ``\emph{multi-agent reinforcement learning autonomous driving}'',
        ``\emph{cooperative driving deep RL}'', ``\emph{vehicle coordination MARL}'',
        and ``\emph{LLM autonomous driving}'';
  \item Forward and backward citation tracing from seed papers
        (Kiran et al.~\cite{kiran2021deep}, Zhang et al.~\cite{zhang2021multi},
        Busoniu et al.~\cite{busoniu2008comprehensive});
  \item Venue coverage: NeurIPS, ICML, ICLR, CVPR, ECCV, ICCV,
        ICRA, IROS, IEEE T-ITS, IEEE T-IV, IEEE T-PAMI, AAAI, IJCAI.
\end{enumerate}

\subsubsection*{Inclusion and Exclusion Criteria}
\textbf{Included}: peer-reviewed conference and journal papers, and
preprints with $\geq 50$ citations that directly address either (a) MARL
methodology applicable to driving, or (b) AD planning/prediction/control
with a learning component.
\textbf{Excluded}: purely perception-only methods without a planning
component, rule-based planners without any learning, and non-vehicular
multi-robot work without demonstrated relevance to road scenarios.
The final corpus comprises \textbf{100 primary references}.

% ------------------------------------------------------------------
\subsection{Three-Layer Taxonomy}
\label{sec:taxonomy}

The organizing framework of this survey consists of three hierarchical
layers that map to successive paradigm shifts in the field:

\begin{enumerate}[leftmargin=*, label=\textbf{Layer \arabic*.}]
  \item \textbf{Algorithmic Foundations (Paradigm I, $\approx$2003--2018).}
        This layer covers the theoretical formalisms (MDP, POMDP, Stochastic
        Games, Dec-POMDP), the deep RL building blocks (DQN, A3C, PPO,
        DDPG, SAC), and the early single-agent driving methods (DAVE-2,
        ChauffeurNet, conditional imitation learning).  The central
        limitation exposed here is that single-agent learning is
        fundamentally insufficient for dense, interactive traffic because
        it cannot account for the non-stationarity induced by other
        learning agents.

  \item \textbf{Coordination and Communication (Paradigm II, $\approx$2017--2022).}
        This layer addresses the core MARL challenge through three parallel
        thrusts: (a) \emph{MARL algorithms} that achieve stable cooperative
        learning (value decomposition, CTDE actor-critic, hierarchical and
        game-theoretic methods); (b) \emph{communication protocols} that
        enable agents to share information selectively and efficiently
        under bandwidth constraints; and (c) \emph{relational representation
        and prediction--planning integration} using GNNs and vectorized
        scene encoders.  The central limitation remaining at the end of
        this paradigm is brittleness in long-tail scenarios that fall
        outside the training distribution.

  \item \textbf{Scalable and Cognitive Architectures (Paradigm III, $\approx$2021--present).}
        This layer covers Transformer-based MARL architectures that scale
        attention over large agent populations, and LLM-driven cognitive
        agents that inject world-scale knowledge and chain-of-thought
        reasoning to handle rare, ambiguous scenarios.  Safety, sim-to-real
        transfer, and world-model-based planning are cross-cutting concerns
        in this layer.
\end{enumerate}

Critically, the layers are not mutually exclusive: modern systems
routinely combine elements from all three (e.g., a GNN scene encoder
from Layer~2 coupled with an LLM reasoning module from Layer~3, trained
via MAPPO from Layer~2, within a CARLA environment from Layer~2).  The
taxonomy is therefore better understood as a set of conceptual lenses
rather than a strict chronological partition.

% ------------------------------------------------------------------
\subsection{Comparison with Existing Surveys and Contributions}
\label{sec:comparison}

A number of surveys have addressed adjacent topics, but none covers the
full scope of this work:

\begin{itemize}[noitemsep]
  \item \textbf{General MARL surveys}~\cite{busoniu2008comprehensive,zhang2021multi}:
        provide excellent algorithmic coverage but lack AD-specific
        application depth, MARL platform comparisons, and the emerging
        LLM/cognitive agent dimension.
  \item \textbf{AD-RL surveys}~\cite{kiran2021deep,sallab2017deep}:
        primarily address single-agent methods; MARL is treated at
        most as an advanced topic in a few paragraphs, with no systematic
        comparison of CTDE methods, communication protocols, or
        value-decomposition approaches.
  \item \textbf{Cooperative driving reviews}: typically narrow in scope,
        covering either traffic-signal control or V2X cooperative
        perception in isolation, without connecting these to the broader
        MARL literature.
  \item \textbf{Trajectory prediction surveys}: excellent coverage of
        motion forecasting methods but generally absent of any connection
        to the policy-learning and multi-agent coordination literature.
\end{itemize}

The unique contributions of this survey are fourfold:
\begin{enumerate}[noitemsep]
  \item \textbf{Unified three-layer taxonomy} that integrates algorithmic,
        representational, and cognitive perspectives under a single
        organizational framework.
  \item \textbf{Systematic comparison of 100 papers}: for each method
        we explicitly state the limitation it addresses, the technical
        mechanism it introduces, and the residual limitations that
        motivated subsequent work.
  \item \textbf{Critical limitation analysis per category}: rather than
        cataloguing results, we trace causal chains between methods,
        making the survey a narrative of scientific progress rather than
        an annotated bibliography.
  \item \textbf{Structured open-challenges mapping}: we distinguish
        between short-term engineering challenges (computational overhead
        of LLM inference, standardizing V2X feature representations)
        and long-term research challenges (formal safety certificates for
        neuro-symbolic systems, generative world models for distributional
        coverage).
\end{enumerate}

% ------------------------------------------------------------------
\subsection{Paper Organization}
\label{sec:organization}

The remainder of this survey is organized as follows.
Section~\ref{sec:background} establishes the formal problem framework,
defining POMDP and POSG formalisms and the AD software stack.
Section~\ref{sec:paradigm1} surveys single-agent learning foundations
(Paradigm~I), covering behavior cloning, covariate shift mitigation,
and single-agent policy optimization.
Section~\ref{sec:paradigm2} is the core of the survey (Paradigm~II),
examining MARL algorithms, communication and coordination protocols,
interaction-aware representation, and autonomous driving applications.
Section~\ref{sec:paradigm3} covers scalable and cognitive architectures
(Paradigm~III), including Transformer-based MARL and LLM-driven agents.
Section~\ref{sec:benchmarks} surveys simulation environments and
evaluation benchmarks.
Section~\ref{sec:challenges} synthesizes open challenges and future
directions.
Section~\ref{sec:conclusion} concludes.

% ==========================================================
\section{Background and Formal Foundations}
\label{sec:background}

This section establishes the mathematical and conceptual foundations upon which
the remainder of the survey is built.  We begin with the single-agent
formalism---Markov Decision Processes (MDPs) and Partially Observable Markov
Decision Processes (POMDPs)---before extending the framework to accommodate
multiple interacting agents through stochastic games and Decentralized POMDPs.
We then briefly review the core deep reinforcement learning algorithms that
underpin the methods discussed in later sections, introduce graph neural networks
as a tool for relational reasoning over traffic scenes, and situate the learning
components within the modular autonomous driving software stack.

% ------------------------------------------------------------------
\subsection{Single-Agent Formalism: MDP and POMDP}
\label{sec:mdp_pomdp}

\subsubsection*{Markov Decision Process}

A \emph{Markov Decision Process} (MDP) provides the canonical mathematical
framework for sequential decision-making under uncertainty.  Formally, an MDP
is defined by the tuple
$\mathcal{M} = \langle \mathcal{S}, \mathcal{A}, P, R, \gamma \rangle$,
where:
\begin{itemize}[noitemsep]
  \item $\mathcal{S}$ is the \textbf{state space}, representing all possible
        configurations of the world (e.g., positions, velocities, and headings
        of all vehicles on the road);
  \item $\mathcal{A}$ is the \textbf{action space}, representing the decisions
        available to the agent (e.g., discrete maneuvers such as accelerate,
        brake, or steer, or continuous values for throttle, steering angle, and
        braking);
  \item $P: \mathcal{S} \times \mathcal{A} \times \mathcal{S} \to [0,1]$ is the
        \textbf{transition probability function}, where $P(s' \mid s, a)$ denotes
        the probability of transitioning to state $s'$ after taking action $a$ in
        state $s$;
  \item $R: \mathcal{S} \times \mathcal{A} \to \mathbb{R}$ is the \textbf{reward
        function}, which encodes the agent's objective (e.g., progress toward the
        goal while penalizing collisions and comfort violations);
  \item $\gamma \in [0, 1)$ is the \textbf{discount factor}, which controls the
        relative weight of immediate versus future rewards.
\end{itemize}
The agent's goal is to find an optimal policy $\pi^*: \mathcal{S} \to \mathcal{A}$
(or a distribution over actions in the stochastic case) that maximizes the
expected discounted cumulative reward
$\mathbb{E}\!\left[\sum_{t=0}^{\infty} \gamma^t R(s_t, a_t)\right]$.
The Markov property---that the next state depends only on the current state and
action, not on the full history---is what makes this optimization tractable via
Bellman recursion.

\subsubsection*{Partially Observable Markov Decision Process}

While the MDP formalism is theoretically elegant, it assumes that the agent
has complete access to the world state $s_t$ at each time step.  This assumption
is violated in virtually every realistic autonomous driving scenario.
Sensor range is finite: a vehicle cannot observe occluded regions behind
buildings or other vehicles.  Sensor resolution is limited: LiDAR point clouds
resolve objects only down to a certain angular precision.  And the world state
includes latent variables that are not directly observable---such as the
\emph{intent} of a pedestrian about to step off the curb, or the
\emph{aggressiveness} of an oncoming merging driver.

The \emph{Partially Observable Markov Decision Process} (POMDP) extends the
MDP to handle this observational uncertainty by adding two components:
\begin{itemize}[noitemsep]
  \item $\mathcal{O}$: an \textbf{observation space}, representing what the
        agent actually perceives (e.g., raw LiDAR scans, bounding-box
        detections, or occupancy grids);
  \item $Z: \mathcal{S} \times \mathcal{A} \times \mathcal{O} \to [0,1]$: an
        \textbf{observation function}, where $Z(o \mid s', a)$ is the probability
        of observing $o$ after transitioning to state $s'$ via action $a$.
\end{itemize}
Because the agent cannot directly observe the true state, the principled
solution is to maintain a \textbf{belief state}
$b_t \in \Delta(\mathcal{S})$---a probability distribution over world states
given the history of observations and actions.  The belief is updated at each
step via Bayes' rule:
\[
  b_{t+1}(s') \;\propto\;
  Z(o_{t+1} \mid s', a_t)\!
  \sum_{s \in \mathcal{S}} P(s' \mid s, a_t)\, b_t(s).
\]
Planning in belief space transforms the POMDP into an MDP over the continuous
belief simplex, which is \textbf{PSPACE-hard} in the finite
horizon~\cite{hubmann2018automated}.  In practice, this intractability is
managed through particle-filter approximations of the belief state, and through
architectures that use recurrent networks (LSTMs or GRUs) or Transformer memory
modules as implicit belief representations, replacing the explicit probability
distribution with a learned latent state.

\citeauthor{hubmann2018automated}~\cite{hubmann2018automated} demonstrate a
concrete instantiation of POMDP planning for autonomous driving at unsignalized
intersections, where the hidden state includes the intent of crossing vehicles.
By planning over a discrete approximation of the belief space, their system can
perform evidence-gathering maneuvers---briefly inching forward to elicit
reactions from occluded drivers---that are impossible for a fully observable MDP
planner.  This example illustrates why the POMDP formalism is not merely a
theoretical refinement but a practically essential one for safety-critical
driving decisions.

% ------------------------------------------------------------------
\subsection{Multi-Agent Extension: Stochastic Games and Dec-POMDP}
\label{sec:multiagent_formalism}

When multiple agents operate simultaneously in a shared environment and each
agent's reward depends on the joint behavior of all agents, the single-agent
POMDP framework becomes insufficient.  Two distinct formalisms are employed
depending on whether agents share objectives (cooperative settings) or have
conflicting incentives (competitive or mixed-motive settings).

\subsubsection*{Stochastic Games and Markov Games}

A \emph{Stochastic Game} (SG), introduced by
\citeauthor{shapley1953stochastic}~\cite{shapley1953stochastic}, is defined by
the tuple
$\langle N, \mathcal{S}, \{\mathcal{A}^i\}_{i=1}^N, P, \{R^i\}_{i=1}^N, \gamma \rangle$,
where $N$ is the number of agents, each with its own action space
$\mathcal{A}^i$ and reward function
$R^i: \mathcal{S} \times \mathcal{A}^1 \times \cdots \times
\mathcal{A}^N \to \mathbb{R}$.
The transition function $P(s' \mid s, a^1, \ldots, a^N)$ depends on the
\emph{joint} action of all agents.  Stochastic games generalize both
single-agent MDPs (when $N=1$) and matrix games (when $|\mathcal{S}|=1$),
and can represent fully competitive (zero-sum), fully cooperative (team-reward),
or mixed-motive multi-agent interactions.

\citeauthor{littman1994markov}~\cite{littman1994markov} introduced the closely
related framework of \emph{Markov Games}, which has become the dominant
formalism in the MARL literature.  Under this framework, the solution concept
is a \emph{Nash equilibrium} policy profile
$(\pi^{1*}, \ldots, \pi^{N*})$ such that no agent can unilaterally improve its
expected return by deviating from its equilibrium policy.  For zero-sum
two-player games, a unique Nash equilibrium in mixed strategies always exists
(guaranteed by minimax theory); for general-sum games with $N > 2$, computing
equilibria is PPAD-hard in the worst case.  In the cooperative case where all
agents share the same reward function ($R^1 = \cdots = R^N = R$), the Markov
Game reduces to a joint policy optimization problem, directly connecting
game-theoretic and MARL formulations.

In the context of autonomous driving, Markov Games capture the strategic
interplay of traffic participants naturally.  A merging scenario, for instance,
can be modeled as a mixed-motive game: the merging vehicle prefers a gap in
traffic, while the through-lane vehicle prefers not to yield---yet both are
penalized for collision.  \citeauthor{sadigh2016planning}~\cite{sadigh2016planning}
exploit this structure to derive trajectories that deliberately influence the
behavior of the human-driven vehicle by communicating intent through motion,
demonstrating how game-theoretic reasoning can make autonomous vehicles more
cooperative and interpretable traffic partners.

\subsubsection*{Decentralized Partially Observable Markov Decision Process}

For cooperative multi-agent settings with partial observability---the dominant
setting in real-world autonomous driving---the most general formalism is the
\emph{Decentralized Partially Observable Markov Decision Process}
(Dec-POMDP)~\cite{bernstein2002complexity}.  A Dec-POMDP augments the Markov
Game with per-agent observation spaces $\{\mathcal{O}^i\}$ and observation
functions $\{Z^i\}$:
\[
  \langle N, \mathcal{S}, \{\mathcal{A}^i\}, P, R, \{\mathcal{O}^i\}, \{Z^i\}, \gamma \rangle.
\]
Crucially, in the fully cooperative Dec-POMDP, agents share a single team
reward $R$, but each agent $i$ observes only
$o^i_t \sim Z^i(\cdot \mid s_t, a_{t-1})$ rather than the full state.
At execution time, each agent must act based solely on its own local
observation history, without access to other agents' observations or actions.

The key theoretical result is that solving a Dec-POMDP optimally is
\textbf{NEXP-hard}~\cite{bernstein2002complexity}---significantly harder than
the PSPACE-hard single-agent POMDP.  The root cause is the combinatorial
explosion in the space of joint histories: with $N$ agents each accumulating
$T$ time steps of observations, the joint observation-action history space
grows as $O\!\left((|\mathcal{O}|\cdot|\mathcal{A}|)^{N \cdot T}\right)$.
In practice, this intractability is addressed through three complementary
strategies:
\begin{enumerate}[noitemsep]
  \item \textbf{CTDE approximations}: training a centralized critic with
        access to global state while restricting each actor to local
        observations only, as formalized by
        \citeauthor{kraemer2016multi}~\cite{kraemer2016multi} and later
        instantiated in MADDPG~\cite{lowe2017multi} and QMIX~\cite{rashid2018qmix};
  \item \textbf{Communication}: allowing agents to share compact messages,
        effectively reducing the degree of partial observability at the cost
        of communication bandwidth (see Section~\ref{sec:paradigm2});
  \item \textbf{Value-function factorization}: assuming conditional
        independence across subsets of agents to decompose the joint
        Q-function into per-agent utilities
        (VDN~\cite{sunehag2018value}, QMIX~\cite{rashid2018qmix}).
\end{enumerate}

A \textbf{cooperative} fleet of autonomous vehicles---such as platoon control
or cooperative intersection management---is most naturally modeled as a
Dec-POMDP: agents share the objective of maximizing traffic throughput while
minimizing total travel time, but each vehicle can only observe agents within
its sensor range and must act locally at execution time.  A
\textbf{mixed-motive} traffic scenario (e.g., lane merging or yielding at a
roundabout) is better modeled as a Markov Game, where individual reward
functions capture private preferences while a shared collision penalty aligns
agents toward safety.

% ------------------------------------------------------------------
\subsection{Deep Reinforcement Learning Primer}
\label{sec:drl_primer}

The methods surveyed in later sections build on a small set of core deep
reinforcement learning algorithms.  This subsection provides a self-contained
summary of these building blocks, emphasizing the aspects most relevant to
their multi-agent extensions.

\subsubsection*{Value-Based Methods: DQN and Variants}

\citeauthor{mnih2015human}~\cite{mnih2015human} demonstrated that a
convolutional neural network trained via Q-learning---the \emph{Deep
Q-Network} (DQN)---could achieve superhuman performance on 49 Atari games
from raw pixel input.  DQN stabilizes Q-learning with two key innovations:
an \emph{experience replay buffer} that breaks temporal correlations between
consecutive training samples, and a \emph{target network} that freezes the
bootstrapping parameters for a fixed number of steps, preventing the
divergence caused by moving learning targets.

In the multi-agent context, DQN forms the foundation of value-decomposition
methods.  Both VDN~\cite{sunehag2018value} and QMIX~\cite{rashid2018qmix}
maintain per-agent local $Q$-functions and mix them into a factored joint
Q-value, inheriting the stability properties of experience replay and target
networks while scaling to cooperative teams of agents.

\subsubsection*{Policy Gradient Methods}

While value-based methods work well for discrete action spaces, autonomous
driving typically demands continuous control (e.g., precise steering angle
and throttle values).  Policy gradient methods directly optimize the
parameters $\theta$ of a stochastic policy $\pi_\theta$ by ascending the
gradient of the expected return:
\[
  \nabla_\theta J(\pi_\theta) = \mathbb{E}_{\pi_\theta}\!\left[
    \nabla_\theta \log \pi_\theta(a_t \mid s_t) \cdot A(s_t, a_t)\right],
\]
where $A(s_t, a_t) = Q(s_t, a_t) - V(s_t)$ is the advantage function that
centers the gradient signal around a baseline, substantially reducing
variance compared to the raw Monte-Carlo return.

\emph{Asynchronous Advantage Actor-Critic} (A3C)~\cite{mnih2016asynchronous}
parallelizes gradient collection across multiple environment instances, enabling
sample-efficient online updates and decorrelating experience.  \emph{Proximal
Policy Optimization} (PPO)~\cite{schulman2017proximal} further stabilizes
training by clipping the policy update ratio within a trust region:
\[
  \mathcal{L}^{\text{CLIP}}(\theta) = \mathbb{E}_t\!\left[
    \min\!\left(r_t(\theta)\, A_t,\;
    \mathrm{clip}(r_t(\theta), 1{-}\epsilon, 1{+}\epsilon)\, A_t\right)\right],
\]
where $r_t(\theta) = \pi_\theta(a_t \mid s_t) / \pi_{\theta_{\text{old}}}(a_t \mid s_t)$
is the probability ratio and $\epsilon$ is the clip threshold.  PPO's
simplicity, stability, and strong empirical performance have made it the
de facto standard for single-agent continuous-control driving tasks, and it
directly generalizes to the multi-agent setting as MAPPO~\cite{yu2021surprising}.

Off-policy methods offer greater sample efficiency.  \emph{Deep Deterministic
Policy Gradient} (DDPG) extends Q-learning to continuous action spaces by
maintaining a deterministic actor $\mu_\theta(s)$ trained by ascending the
critic gradient $\nabla_\theta Q(s, \mu_\theta(s))$.  \emph{Soft Actor-Critic}
(SAC) adds entropy regularization to the policy objective, encouraging
exploration and preventing premature convergence to deterministic policies---a
property particularly beneficial in sparse-reward driving scenarios where
broad exploration of the state space is critical.

\subsubsection*{Relevance to Multi-Agent Extensions}

Each single-agent algorithm above has a multi-agent analogue:
DQN $\to$ QMIX~\cite{rashid2018qmix} (cooperative value decomposition);
DDPG $\to$ MADDPG~\cite{lowe2017multi} (centralized critic, decentralized
actors); PPO $\to$ MAPPO~\cite{yu2021surprising} (agents with a shared or
centralized critic).  The key challenge in adapting single-agent algorithms
to the multi-agent setting is the \emph{non-stationarity} problem: from
agent $i$'s perspective, the environment's transition distribution
$P(s' \mid s, a^i)$ appears non-stationary because it implicitly depends on
the simultaneously changing policies of all other agents.  This violates the
Markov assumption required for convergence guarantees of standard single-agent
RL algorithms, and motivates the Centralized Training, Decentralized Execution
(CTDE) framework discussed at depth in Section~\ref{sec:paradigm2}.

% ------------------------------------------------------------------
\subsection{Graph Neural Networks for Relational Reasoning}
\label{sec:gnns}

Traffic scenes are inherently \emph{relational}: the relevance of a
neighboring vehicle to the ego-agent depends on spatial proximity, relative
heading, and the road topology.  Standard convolutional neural networks
process regular grids and lack the inductive biases needed to efficiently
encode these pairwise agent relationships.  \emph{Graph Neural Networks}
(GNNs) provide the appropriate representation: they model a traffic scene as
a directed graph $\mathcal{G} = (\mathcal{V}, \mathcal{E})$ where nodes
$v \in \mathcal{V}$ represent agents (or road elements) and edges
$(u, v) \in \mathcal{E}$ encode spatial or semantic interactions.
Node features are iteratively updated via message passing:
\[
  h_v^{(k+1)} = \mathrm{UPDATE}\!\left(h_v^{(k)},\;
    \mathrm{AGGREGATE}\!\left(\{m_{uv}^{(k)} : (u,v) \in \mathcal{E}\}\right)\right),
\]
where $m_{uv}^{(k)}$ is a message from neighbor $u$ to $v$ at layer $k$.

\emph{Graph Attention Networks} (GAT)~\cite{velickovic2018graph} extend this
framework by computing learned attention coefficients $\alpha_{uv}$ that
weight each neighbor's contribution:
\[
  \alpha_{uv} = \frac{
    \exp\!\left(\mathrm{LeakyReLU}(\mathbf{a}^\top
      [\mathbf{W} h_u \,\|\, \mathbf{W} h_v])\right)
  }{
    \sum_{k \in \mathcal{N}(v)} \exp\!\left(\mathrm{LeakyReLU}(\mathbf{a}^\top
      [\mathbf{W} h_k \,\|\, \mathbf{W} h_v])\right)
  },
\]
where $\mathbf{W}$ is a learnable projection matrix, $\mathbf{a}$ is an
attention vector, and $\|$ denotes vector concatenation.  The attention
mechanism allows the network to focus automatically on the most relevant
neighboring agents---such as the lead vehicle in the ego-lane during a
car-following maneuver---without requiring hand-crafted distance thresholds
or fixed neighborhood sizes.

\citeauthor{jiang2019graph}~\cite{jiang2019graph} apply GNNs directly to
cooperative MARL, constructing a dynamic graph from agent observations and
using graph convolutions to compute communication-like feature aggregations.
Their Graph Convolutional Reinforcement Learning approach demonstrates that
policies learned over a graph adjacency structure are naturally invariant to
permutations of agent indices---a critical inductive bias when the number of
surrounding vehicles changes dynamically across episodes.

In the trajectory prediction domain, \citeauthor{salzmann2020trajectron}
\cite{salzmann2020trajectron} (Trajectron++) and
\citeauthor{liang2020lanegcn}~\cite{liang2020lanegcn} (LaneGCN) demonstrate
that GNNs encoding heterogeneous agent-to-agent and agent-to-map interactions
substantially outperform architectures based on rasterized top-down images,
both in displacement error and in the calibration of predicted probability
distributions.

This subsection provides the representational foundation for the
value-decomposition and communication methods reviewed in
Section~\ref{sec:paradigm2}: QMIX, graph-based MARL, and VectorNet all rely
on variants of graph-structured feature aggregation to encode inter-agent
dependencies compactly.

% ------------------------------------------------------------------
\subsection{Overview of the Autonomous Driving Software Stack}
\label{sec:ad_stack}

Understanding where learning-based methods intervene in the broader autonomous
driving pipeline is essential for contextualizing the papers reviewed in
subsequent sections.  Classical autonomy systems adopt a \emph{modular
pipeline} architecture consisting of four sequentially composed sub-systems:

\begin{enumerate}[leftmargin=*, label=\textbf{(\arabic*)}]
  \item \textbf{Perception.}
        Raw sensor data (LiDAR point clouds, camera images, radar returns)
        are processed to detect and classify objects (vehicles, pedestrians,
        cyclists), estimate their 3-D poses, and build a semantic map of the
        scene.  While deep learning has largely solved static-object detection,
        multi-sensor fusion under adverse weather and detection under adversarial
        conditions remain active research areas.

  \item \textbf{Prediction.}
        Given the current perceived state, the prediction module forecasts
        the short-horizon ($\leq$5 s) future trajectories of all road users.
        This is the module most directly connected to the interaction-modeling
        literature surveyed in Section~\ref{sec:paradigm2}: Social
        GAN~\cite{gupta2018social}, Trajectron++~\cite{salzmann2020trajectron},
        VectorNet~\cite{gao2020vectornet}, and TNT~\cite{zhao2020tnt} all
        operate at this layer.

  \item \textbf{Planning.}
        Given predicted futures, the planner selects a reference trajectory
        or sequence of waypoints for the ego-vehicle that optimizes the driving
        objective (efficiency, comfort, safety) while respecting traffic rules
        and predicted agent behaviors.  MARL methods primarily intervene at this
        layer: cooperative policy networks replace or augment the traditional
        motion planner, enabling dynamic negotiation of right-of-way and
        adaptive formation control.

  \item \textbf{Control.}
        A low-level controller (PID, MPC, or neural feedback policy) tracks the
        planned trajectory by issuing throttle, braking, and steering commands
        to the vehicle actuators.
\end{enumerate}

The modular pipeline's principal weakness is the \emph{cascading error}
problem: imperfections at any stage compound as they propagate downstream, and
the planning module receives no gradient signal that would allow it to request
more uncertainty-aware predictions.  Furthermore, the pipeline's clean
separation of concerns breaks down whenever prediction and planning are tightly
coupled, as in the merged prediction-planning paradigm of PiP~\cite{song2020pip}.

\emph{End-to-end} (E2E) learning approaches collapse the entire pipeline into
a single neural network that maps raw observations directly to control
commands~\cite{bojarski2016end,codevilla2018end}.  E2E architectures eliminate
the information bottleneck between modules and allow the system to learn
task-relevant representations directly from driving objectives.  However, E2E
methods trade interpretability for performance and raise significant challenges
for safety certification.  \citeauthor{sallab2017deep}~\cite{sallab2017deep}
propose a recurrent deep RL architecture positioned between these extremes:
it uses learned perceptual features but trains a separate policy module via RL
rather than pure imitation, gaining flexibility without discarding modularity
entirely.

MARL methods can be positioned along this modular-to-end-to-end spectrum.
Value-decomposition approaches such as QMIX~\cite{rashid2018qmix} typically
operate at the planning level, receiving pre-processed observation vectors
from an upstream perception stack.  Communication-based methods like
TarMAC~\cite{das2019tarmac} add a lateral information-sharing channel between
the planning modules of multiple vehicles.  LLM-based cognitive agents
(Section~\ref{sec:paradigm3}) operate at the highest level of abstraction,
producing natural-language driving instructions that are translated into motion
primitives by downstream controllers.

Together, the formal frameworks (MDP/POMDP/SG/Dec-POMDP), algorithmic
building blocks (DQN, PPO, DDPG, actor-critic), representational tools (GNNs),
and architectural context (the AD software stack) established in this section
provide the essential scaffolding for the three-paradigm survey that follows.

% ==========================================================
\section{Paradigm I: Single-Agent Learning Foundations}
\label{sec:paradigm1}

The first paradigm in the evolution of learning-based autonomous driving
(approximately 2016--2018) established the core machine-learning machinery
that all subsequent multi-agent methods build upon.  During this period,
researchers demonstrated that deep neural networks could learn end-to-end
driving policies directly from raw sensory input, that reinforcement learning
could optimize policies in simulated traffic without hand-engineered reward
functions, and that interaction-aware trajectory prediction models could
forecast the short-horizon futures of neighboring road users.  These
contributions were transformative.  Yet they shared a common, fundamental
blind spot: each method treated the traffic environment as if it contained
only a single learning agent, ignoring the strategic interdependence between
co-evolving policies.  This section surveys the key methods of Paradigm~I,
assessing both their technical contributions and the precise limitations that
motivated the multi-agent methods of Paradigm~II.

% ------------------------------------------------------------------
\subsection{Imitation Learning and Behavior Cloning}
\label{sec:imitation}

\subsubsection{End-to-End Learning: DAVE-2 and Successors}

The dominant paradigm of the pre-deep-learning era decomposed driving into a
pipeline of hand-engineered modules.  The landmark work of
\citeauthor{bojarski2016end}~\cite{bojarski2016end} (the NVIDIA DAVE-2 system)
challenged this paradigm with a radical simplification: a nine-layer
convolutional neural network trained end-to-end to map raw front-facing camera
images directly to a continuous steering command.  Trained on approximately
72 hours of human driving data, DAVE-2 successfully navigated roads with
explicit lane markings, roads without lane markings, and even parking lots in
experiments, demonstrating that a single feed-forward network could implicitly
learn the full perception-to-control mapping without any intermediate
supervision.

The contribution of DAVE-2 was primarily proof of concept: it established that
deep imitation learning for driving was feasible at all.  Three concrete
technical advantages emerged from the end-to-end formulation.  First, the
network could learn task-relevant features automatically, without the need to
specify which visual cues correspond to lane boundaries or upcoming obstacles.
Second, training was data-efficient in the sense that the same labeled examples
simultaneously improved all components of the pipeline, rather than requiring
separate annotation budgets for perception, prediction, and planning.  Third,
the absence of hand-crafted intermediate representations eliminated the
information bottleneck inherent in multi-stage pipelines.

\citeauthor{codevilla2018end}~\cite{codevilla2018end} extended this paradigm
with \emph{Conditional Imitation Learning} (CIL), addressing the crucial
limitation that a pure regression from image to steering cannot resolve the
ambiguity of navigational goals.  At an intersection, for instance, the optimal
steering command is entirely determined by whether the destination lies ahead,
to the left, or to the right---information absent from any single camera frame.
CIL conditions the policy on a high-level navigational command (``turn left'',
``go straight'', ``turn right'', ``follow road'') provided by a route planner,
using a branched network architecture where each branch is activated only when
the corresponding command is issued.  On the CARLA simulator benchmark, CIL
substantially outperformed unconditional imitation baselines, establishing the
conditional architecture as the standard recipe for end-to-end urban driving.

Despite these advances, end-to-end imitation learning suffers from a
structural weakness known as \emph{covariate shift}~\cite{bansal2018chauffeur}.
A behavior-cloned policy is trained exclusively on states visited by the expert
demonstrator, but at test time the agent's own imperfect actions progressively
deviate from the expert's trajectory, leading the agent into novel states
$s \notin \mathcal{D}_{\text{train}}$ for which the cloned policy has no
reliable prediction.  Small errors compound: a slight drift to the right of the
lane at time $t$ places the agent in an out-of-distribution state at time
$t+1$, and the policy---never trained to recover from such states---may respond
with an erratic command that amplifies the error further.  In the single-agent
setting, this is already a critical failure mode; in multi-agent traffic, where
the state space is enlarged by the behaviors of surrounding vehicles, the
covariate shift problem becomes dramatically more severe.

\subsubsection{Addressing Covariate Shift: DAgger and Data Augmentation}

The theoretical solution to covariate shift is \emph{Dataset Aggregation}
(DAgger), introduced by \citeauthor{ross2011reduction}~\cite{ross2011reduction}.
DAgger is an iterative on-policy imitation learning algorithm: at each round,
the current learned policy $\hat{\pi}_t$ is rolled out in the environment,
collecting a dataset of states $\mathcal{D}_t$ actually visited by the
(imperfect) learner.  Expert labels are then queried for those visited states,
and the policy is retrained on the growing aggregate dataset
$\mathcal{D}_1 \cup \cdots \cup \mathcal{D}_t$.  By ensuring that the
training distribution matches the \emph{learner's} state distribution rather
than the expert's, DAgger provides a no-regret guarantee: the cumulative
cost of the learned policy converges to that of the best policy in the
hypothesis class.  The catch is that DAgger requires an interactive expert
capable of labeling any state the learner might visit---expensive in terms of
human annotation or oracle computation.

\citeauthor{bansal2018chauffeur}~\cite{bansal2018chauffeur} addressed the same
problem at industrial scale in \emph{ChauffeurNet}, trained on over 30 million
labeled driving examples from Waymo's real-world fleet.  Their key finding was
counterintuitive: even at this scale, a policy trained on clean expert
demonstrations failed to generalize to novel test scenarios.  Recovery from
suboptimal states---a near-collision, a vehicle temporarily off the lane
center---could not be learned from demonstrations that never showed such
situations.  The solution was \emph{perturbation-based data augmentation}: the
training data was synthetically corrupted by injecting random perturbations into
the rendered bird's-eye-view inputs, and the expert labels were queried for
these perturbed states.  ChauffeurNet operates on a mid-level representation
(a top-down rendering of the scene including agent boxes, road structure, and
traffic lights) rather than raw pixels, making it possible to generate
perturbed variants analytically.  The resulting policy demonstrated markedly
improved robustness to distributional shift and was shown to transfer to
real-vehicle trials.

A critical limitation of ChauffeurNet and DAgger-style approaches persists,
however: the perturbations are applied in isolation, treating the ego-vehicle
as the only actor whose state can deviate from nominal.  In reality, the
distribution of \emph{other} vehicles' states---the speed of the car ahead
that is braking unexpectedly, the trajectory of a pedestrian who begins to
cross at an atypical moment---is equally unconstrained at test time.  No
synthetic augmentation scheme that treats only the ego-vehicle as a random
variable can capture the full diversity of multi-agent traffic interaction.
The fundamental gap is architectural: behavior cloning models a policy
$\pi(a \mid o)$ that responds to observations of other agents but does not
\emph{reason about} them as co-equal agents pursuing their own objectives.

% ------------------------------------------------------------------
\subsection{Reinforcement Learning for Single-Agent Driving}
\label{sec:rl_single}

Reinforcement learning removes the requirement for expert demonstrations by
optimizing the policy directly against a reward signal.  The agent improves
through trial-and-error interaction with the environment, discovering
behaviors that maximize cumulative reward without supervision from
demonstrations.

\citeauthor{sallab2017deep}~\cite{sallab2017deep} proposed one of the earliest
applications of deep RL to autonomous driving, combining a recurrent
architecture (LSTM-augmented policy network) with an asynchronous actor-critic
training scheme.  The recurrent state allowed the policy to maintain an
implicit belief over the full driving history, partially compensating for the
limited field of view of a single camera.  Their framework was evaluated in a
lane-keeping and lane-change scenario in a simulated highway environment,
demonstrating that policies trained with deep RL could match and exceed the
performance of rule-based baselines without any explicit traffic rules being
programmed.

For continuous-action driving tasks---where the output is a real-valued
steering angle and throttle rather than a discrete maneuver selection---
Proximal Policy Optimization~\cite{schulman2017proximal} emerged as the
dominant training algorithm due to its balance of sample efficiency, stability,
and ease of implementation.  Unlike TRPO, which solves a constrained
optimization at each step, PPO enforces the trust-region constraint
approximately via a clipped surrogate objective, making it computationally
tractable for long-horizon driving rollouts.  PPO-trained policies for
lane-following, intersection navigation, and roundabout negotiation in CARLA
became the standard single-agent baselines against which multi-agent methods
are later evaluated.

\citeauthor{kahn2017uncertainty}~\cite{kahn2017uncertainty} addressed an
orthogonal limitation of standard deep RL for driving: the inability to reason
about uncertainty in the learned value function.  Standard model-free RL
explores aggressively, which is acceptable in video-game environments but
catastrophic in physical or realistically simulated driving.  Their
\emph{PLATO} system maintained an ensemble of neural network dynamics models,
using the disagreement between ensemble members as a proxy for epistemic
uncertainty.  A safety-constrained planner avoided actions that led to
high-uncertainty states, dramatically reducing collision rates during the
exploration phase.  This work anticipated the broader \emph{safe RL}
literature that became crucial once MARL methods were deployed in real traffic.

\citeauthor{chen2019model}~\cite{chen2019model} tackled the structural
complexity of urban driving by decomposing the policy into hierarchical
sub-tasks.  Their model-free hierarchical RL framework trained a high-level
meta-controller to select among driving primitives (``change to left lane'',
``follow current lane'', ``decelerate'') and a set of low-level primitive
controllers to execute each selected behavior.  This temporal abstraction
reduced the effective planning horizon for each sub-policy and allowed
different primitives to be trained in targeted scenarios rather than requiring
all behaviors to emerge from a single flat policy.  The approach achieved
substantially better performance on urban intersection scenarios with
uncontrolled traffic in CARLA compared to flat PPO baselines, at the cost of
requiring manual specification of the primitive taxonomy.

Despite these advances, single-agent RL methods share a critical structural
limitation when deployed in multi-vehicle environments: from any individual
agent's perspective, other vehicles are part of the environment, and if those
vehicles are also learning agents (as in any MARL training run), the transition
function $P(s' \mid s, a)$ experienced by the ego-agent is \emph{non-stationary}.
As the policies of surrounding vehicles update during training, the
distribution of their actions---and hence the joint transition distribution---
shifts continuously.  This violates the stationarity assumption underlying
all convergence guarantees for single-agent RL algorithms.  The empirical
consequence is well-documented: single-agent PPO policies that perform
excellently in scenarios with scripted, rule-based traffic participants
degrade significantly when evaluated against other trained-RL agents, because
the behavioral distribution of those agents was never seen during training.

% ------------------------------------------------------------------
\subsection{Interaction-Aware Trajectory Prediction (Single-Agent Perspective)}
\label{sec:prediction_single}

In parallel with the policy-learning literature, a distinct research thread
focused on improving the \emph{prediction} of other agents' future trajectories
from the perspective of a single observing ego-vehicle.  Accurate trajectory
prediction is a prerequisite for safe planning: the ego-agent needs short-horizon
forecasts of surrounding agents to evaluate whether a given maneuver (e.g.,
a lane change) will produce a collision.

The seminal work in this direction is Social LSTM by
\citeauthor{alahi2016social}~\cite{alahi2016social}, which modeled pedestrian
trajectory prediction in crowded spaces.  A key contribution was the
\emph{Social Pooling} mechanism: rather than predicting each pedestrian's
trajectory independently, the LSTM hidden states of all pedestrians in a
spatial neighborhood were pooled together into a joint context vector that
informed each individual's prediction.  This enabled the model to capture
interaction-induced behaviors such as collision avoidance and group formation.
In the driving context, the grid-based pooling was adapted to model how
neighboring vehicles' current velocities and headings influence the ego-vehicle's
trajectory distribution.  \citeauthor{deo2018convolutional}~\cite{deo2018convolutional}
refined this approach with \emph{Convolutional Social Pooling}, replacing the
fixed-grid occupancy representation with a learned convolutional aggregation
that better captured spatial structure in highway merging scenarios.

\citeauthor{lee2017desire}~\cite{lee2017desire} (DESIRE) introduced a
more powerful generative modeling framework for trajectory prediction,
combining a conditional variational autoencoder (CVAE) with a recurrent
inverse RL component.  DESIRE generates a diverse set of candidate future
trajectories using the CVAE's latent code and then ranks them with a
learned scoring function that evaluates social acceptability and adherence
to the inferred reward function.  The key advance over Social LSTM was the
\emph{multi-modal} output: rather than predicting a single most-likely
trajectory, DESIRE explicitly samples from a distribution of plausible
futures, which is essential when agents face genuine decision ambiguity
(e.g., whether to merge before or after an approaching vehicle).

\citeauthor{gupta2018social}~\cite{gupta2018social} (Social GAN) pushed the
multi-modal paradigm further by replacing the CVAE with a generative
adversarial network (GAN).  The generator samples diverse trajectory
proposals conditioned on a \emph{Global Pooling} mechanism that aggregates
all agents' encoded histories into a single scene-level feature vector.
A discriminator is trained to distinguish generated trajectories from
ground-truth ones, enforcing social plausibility of the output distribution.
Social GAN was among the first works to demonstrate that diversity of
predicted trajectories could be evaluated quantitatively via metrics such as
the \emph{Minimum over $K$ samples} displacement error (minFDE-$K$), which
evaluates how close the best of $K$ generated trajectories is to the
ground truth.  This metric subsequently became a standard benchmark
across the prediction literature.

A critical limitation shared by Social LSTM, DESIRE, and Social GAN is the
\emph{prediction-planning disconnect}: all three systems predict the
trajectories of surrounding agents assuming the \emph{ego-agent's trajectory
is fixed}.  In reality, the ego-vehicle's own planned action influences
what neighboring agents will do.  A vehicle planning to brake suddenly
elicits very different behaviors from the vehicle behind it than a vehicle
planning to accelerate.  This bidirectional coupling---prediction depends on
planning, planning depends on prediction---cannot be captured by any model
that treats prediction as a feed-forward module upstream of a separate planner.
The merged prediction-planning paradigm (PiP~\cite{song2020pip}) that addresses
this coupling appears only in Paradigm~II, once the representational tools
(VectorNet, graph neural networks) for encoding it compactly became available.

% ------------------------------------------------------------------
\subsection{Critical Analysis: Limitations Driving Paradigm~II}
\label{sec:paradigm1_limits}

The methods surveyed in this section collectively constitute Paradigm~I:
a set of powerful tools, each optimal within its problem formulation, but
each resting on a single-agent assumption that breaks down systematically
in dense, interactive traffic.  Three concrete failure modes emerged clearly
from the literature and directly motivated the MARL methods of Paradigm~II.

\textbf{(1) Covariate shift in imitation learning.}
Behavior-cloned policies are trained on expert state distributions but
evaluated on their own (drifting) state distributions.  DAgger and
perturbation-based augmentation mitigate but do not eliminate this gap;
crucially, neither addresses the co-evolution of surrounding agents'
states.  The structural fix---learning from a reward signal via RL rather
than pure imitation, or from diverse multi-agent interaction data---is the
entry point for Paradigm~II's policy-gradient and CTDE methods.

\textbf{(2) Non-stationarity in single-agent RL.}
Deploying a single-agent RL policy in a multi-vehicle environment makes
the transition function non-stationary from the ego-agent's perspective.
Convergence guarantees evaporate, and policies trained against scripted
opponents fail to generalize to adaptive opponents.  The CTDE framework
(Section~\ref{sec:paradigm2})---in which a centralized critic conditions
on the joint action of all agents during training while decentralized
actors act locally at execution---resolves non-stationarity at training
time by explicitly acknowledging the multi-agent structure of the value
function.

\textbf{(3) Prediction-planning disconnect.}
Single-agent perspective trajectory predictors model other agents as
external stochastic processes, ignoring the ego-vehicle's reactive
influence on those processes.  This is acceptable for distant agents
or low-traffic-density scenarios, but fails catastrophically in
tight merging and yielding interactions where agents continuously
update their intentions in response to each other.  The bidirectional
prediction-planning methods of Paradigm~II (PiP~\cite{song2020pip},
game-theoretic planners~\cite{sadigh2016planning,fisac2019hierarchical})
address this coupling directly.

Together, these three failure modes define the problem that Paradigm~II
sets out to solve: endowing autonomous vehicles with a genuinely
\emph{multi-agent} view of traffic, in which all participants are treated
as co-equal agents with evolving policies rather than as independently
sampled stochastic obstacles.

% ==========================================================
\bibliographystyle{unsrtnat}
\bibliography{references}

\end{document}
